% Figures moved from the main paper to shorten it; regenerated by run_all_paper_experiments.py.

\begin{figure}[!ht]
	\centering
	\includegraphics[width=\textwidth]{fig07_error_cdf.pdf}
	\caption{Empirical cumulative distributions of the positioning error of RF300-F1.0 and RF300-F0.7 on the official Building, Office, and Apartment test sets.}
	\label{fig:s_error_cdf}
\end{figure}

\begin{figure}[!ht]
	\centering
	\includegraphics[width=\textwidth]{fig08_grouped_cv.pdf}
	\caption{Fold-level results of five-fold RP-grouped validation on the original benchmark (pooled values in the RP-grouped cross-validation table of the paper).}
	\label{fig:s_grouped_cv}
\end{figure}

\begin{figure}[!ht]
	\centering
	\includegraphics[width=0.8\textwidth]{fig09_eetac_external_validation.pdf}
	\caption{External validation on the EETAC auditorium dataset (frame-aligned, seed 42) using sample-level and RP-grouped protocols.}
	\label{fig:s_eetac}
\end{figure}

\begin{figure}[!ht]
	\centering
	\includegraphics[width=0.8\textwidth]{fig11_cross_day.pdf}
	\caption{Cross-day EETAC results after frame alignment (seed 42). Each model is trained on one acquisition day and evaluated on the other.}
	\label{fig:s_cross_day}
\end{figure}

\begin{figure}[!ht]
	\centering
	\includegraphics[width=0.8\textwidth]{fig06_original_accuracy.pdf}
	\caption{Paper-compatible error on the original benchmark: published hybrid RF, reproduced RF300-F1.0, and RF300-F0.7 (official RP-disjoint test sets).}
	\label{fig:s_original_accuracy}
\end{figure}

\begin{figure}[!ht]
	\centering
	\includegraphics[width=0.8\textwidth]{fig10_accuracy_efficiency.pdf}
	\caption{Accuracy versus inference time of the 300-tree RF300-F1.0 and RF300-F0.7 models (values in Table~\ref{tab:s_efficiency}).}
	\label{fig:s_accuracy_efficiency}
\end{figure}

\begin{figure}[!ht]
	\centering
	\includegraphics[width=0.8\textwidth]{fig12_uncertainty_area_coverage.pdf}
	\caption{Prediction-region area versus empirical coverage for circular and covariance-aware elliptical regions in a single calibration split (seed 42; values in Table~\ref{tab:s_uncertainty}).}
	\label{fig:s_uncertainty_area}
\end{figure}

\begin{figure}[!ht]
	\centering
	\includegraphics[width=\textwidth]{fig14_ratio_selection.pdf}
	\caption{Error versus subspace ratio as seen by the official RP-disjoint test, five-fold RP-grouped CV on the training RPs, and the out-of-bag error. The dashed line marks the fixed default $\rho=0.7$ (regret of each selection rule in the ratio-selection table of the paper).}
	\label{fig:s_ratio_selection}
\end{figure}
